\section{Mistral 7B：data-first 与紧凑模型设计}

Mistral 7B 的课堂故事延续了两个判断：第一，数据质量是小团队最值得投入的杠杆；第二，小模型只有同时优化训练和推理架构，才能把“参数少”转化为真实部署优势。Mistral 7B 官方论文给出两个关键机制：Grouped-Query Attention（GQA）和 Sliding Window Attention（SWA）。

\subsection{GQA 与 SWA 解决什么问题}

\begin{knowledgebox}{术语消化：紧凑模型的注意力机制}
\begin{tabularx}{\textwidth}{L{0.18\textwidth}>{\raggedright\arraybackslash}X>{\raggedright\arraybackslash}X}
\toprule
术语 & 解决的问题 & 核心机制与部署关系 \\
\midrule
GQA & KV cache 和解码带宽过大 & 多个 query heads 共享较少的 key/value heads，降低 KV cache 与读取量，提升 inference throughput。 \\
SWA & 长上下文 full attention 的二次复杂度 & 每层只直接关注局部窗口，信息通过多层逐步传播，以较低成本扩大有效感受野。 \\
KV cache & 自回归解码重复计算历史 key/value & 保存历史注意力状态以避免重算；其容量常成为 batch size 和并发的主要约束。 \\
\bottomrule
\end{tabularx}
\end{knowledgebox}

Mistral 7B 论文还以 Apache 2.0 发布权重和参考实现。这里的系统意义不是“许可证自动创造生态”，而是允许团队在本地或云上部署、修改和 fine-tune，从而把基础模型当作可组合资产。开放权重扩大了 experimentation surface，也把安全、评测和运维责任更多地交给集成者。

\begin{importantbox}{小模型优势必须落到四个资源维度}
\begin{enumerate}
  \item \textbf{显存}：参数、KV cache 和 runtime workspace 能否放入目标硬件；
  \item \textbf{带宽}：每个 token 需要从 HBM 读取多少权重和 cache；
  \item \textbf{延迟}：用户是否能接受首 token 与逐 token 延迟；
  \item \textbf{吞吐}：在给定服务等级下能容纳多少并发请求。
\end{enumerate}
Benchmark 领先若不能改善这些资源账，就未必形成部署优势。
\end{importantbox}

\begin{warningbox}{架构论文不能替代 workload benchmark}
GQA 和 SWA 提供机制层优势，但真实收益受序列长度、batching、kernel 实现、量化、硬件和服务目标影响。选择模型时应在目标流量与数据上测量，而不是只复制论文平均分。
\end{warningbox}

\subsection{本章小结}

Mistral 7B 的可教学价值在于把 data-first 组织方式与 inference-aware 架构连接起来。紧凑模型不是缩水版大模型，而是围绕可部署资源约束重新分配训练与架构预算。

